\subsection{李雅普诺夫式推理}

\begin{definition}[李雅普诺夫式判据]
\textbf{李雅普诺夫式判据}是一个沿轨迹持续下降的标量量，
它因此为系统正在趋向某个稳定集合提供证据。
\end{definition}

本书引入李雅普诺夫语言，并不是为了对日常常规执行严格证明。
它的目的在于引入一种设计理想：
干预是否能够提供一个单调变化的进展指标，而不是只依赖含混的自我报告？

\begin{example}[单调进展代理指标]
假设一位读者想建立一个阅读机制。
一个有用的实践代理指标，不应是“我今天有没有感觉自己很学术”，
而应是“我是否在固定时间之前写下了第一条页边批注”。
这不是一个数学上的李雅普诺夫证明。
但它是一种李雅普诺夫式设计动作：选择一个可测量量，其持续改善能够支撑如下假设，即目标机制正在变得更容易进入。
\end{example}